% Compared with arXiv:2512.15383v1; AUTHOR-CHECK notes remain author decisions.
\documentclass[10pt]{article}
\usepackage[margin=0.72in]{geometry}
\usepackage[T1]{fontenc}
\usepackage{lmodern,amsmath,enumitem,xcolor,microtype}
\usepackage{xr-hyper}
\usepackage[colorlinks=true,linkcolor=black,urlcolor=blue,bookmarks=false]{hyperref}
\externaldocument[M-]{main}[main.pdf]
\setlength{\parindent}{0pt}
\setlength{\parskip}{0.35em}
\setlength{\emergencystretch}{2em}
\setlist[itemize]{leftmargin=1.25em,itemsep=0.55em,topsep=0.3em}
\newcommand{\mssec}[1]{Section~\ref*{M-#1}}
\newcommand{\msapp}[1]{Appendix~\ref*{M-#1}}
\newcommand{\checkauthor}[2]{\textcolor{red!65!black}{\textbf{AUTHOR-CHECK #1:} #2}}
\begin{document}
\begin{center}
{\Large\bfseries Summary of Revisions}\par
\vspace{0.3em}
{\bfseries Interpretable Multivariate Conformal Prediction\par
with Fast Transductive Standardization}\par
JMLR-25-3153-1
\end{center}

This summary compares the current draft with the original arXiv v1
submission of December 17, 2025. It describes changes present in the draft;
red author checks identify unresolved decisions rather than completed changes.

\begin{itemize}
\item \textbf{Abstract.} The central contribution and finite-sample coverage
claim are retained. Grammatical corrections improve readability; no new
theoretical result is claimed.

\item \textbf{Section 1: introduction and related work.} The motivation for
interpretable simultaneous intervals and the Bonferroni comparison are
retained. Figure 1 now shows illustrative marginal coverages of 92\% and
93\% alongside 90\% joint coverage, making that distinction visible
immediately. The final pass tightens the language.
\checkauthor{LIT}{The Point CHR/CQHR distinction, shape-template discussion,
and empirical attribution still need substantive review.}

\item \textbf{Section 2: setup and motivating constructions.} The revised
setup states joint exchangeability of calibration and test observations
and explains that coverage is marginal over the data but simultaneous
across outputs. Residual definitions discuss capping and ties. A common
location/scale parameterization separates the population oracle from its
practical approximations. These changes clarify the statistical target and
motivation for coordinate adaptation; vector/scalar notation is standardized.
\checkauthor{MAP/TIES/BALANCE}{Review inverse terminology, tie handling, and
the coverage-balance claim.}

\item \textbf{Section 3: methodology.} A threshold-indexed set map connects
the transductive oracle, link functions, and global/local bounds. The local
argument now gives a pointwise implication on oracle acceptance and
working-region membership, not cell-conditional coverage. Figure 2 labels
the residual boundaries; multi-indices are consistently bold. The
real-data diagnostic reference now points to the appendix. These changes
make the construction easier to follow.
\checkauthor{EVENT/ENDPOINT/BOUNDARY}{The piecewise proof, search-lemma
endpoint, and partition boundaries still need mathematical review.}

\item \textbf{Section 4: numerical experiments.}
\textbf{Twelve experimental studies have been added in total: ten simulation
studies and two real-data diagnostic studies.} This count spans Section 4
and Appendix C and treats a parameter sweep as one study. Additional
panels and coordinate plots are not counted separately. E11 and E12 share
a rerun cohort, and E01 and E04 use the same dependent-Gaussian experiment
family; the count is of distinct study questions, not independent datasets.
The original independent Gaussian benchmarks, six-dataset comparison,
and energy illustration remain.
\item \textbf{E01: Dependent heterogeneous Gaussian noise (\mssec{sec:absolute-residual-experiments}).}
Tests whether coordinate adaptation remains useful when outputs are correlated. Joint coverage, volume, runtime, and coordinate lengths distinguish scale adaptation from the smaller but under-covering copula regions.

\item \textbf{E02: Partial heteroskedasticity (\mssec{sec:absolute-residual-experiments}).}
Makes only five of ten coordinates input-dependent while preserving marginal second moments. Separates heteroskedasticity from unequal marginal scales and assesses joint, not input-conditional, coverage.

\item \textbf{E03: CQR with native CQHR (\mssec{sec:cqr-experiments}).}
Tests quantile-regression residuals and adds the requested CQHR baseline. Outcome-space volumes and lengths show design-dependent gains of capped TSCP while identifying its inability to shrink the fitted base interval.

\item \textbf{Section 4.4: real-data presentation.} The original results
are retained in an expanded table. A new plot of existing timing records
makes wall-clock cost visible in the body; it is not counted as an added
experiment. The Point CHR infinite volume on \texttt{stock} is explained
by the small final calibration split. Coordinate and search diagnostics
are presented in Appendix C.

\item \textbf{Section 5: discussion.} The existing future directions are
retained. New cross-references connect the outlier discussion to the
contamination study and distinguish a direct signed-score extension from
the capped and shifted CQR experiments.
\end{itemize}

\clearpage
\begin{itemize}
\item \textbf{Appendix A: algorithmic details and complexity.} The algorithms
and complexity bounds are retained, with standardized vector notation and
grammar. Branch diagnostics complement rather than improve the worst-case
bound. \checkauthor{DS}{Review the final-subset indexing in the data-splitting
algorithm's quantile.}

\item \textbf{Appendix B: proofs.} The original arguments are retained with
language, notation, punctuation, and display-layout corrections. No new
theorem for infinite variance, ties, or signed residuals is claimed.
\checkauthor{DOMAIN/ZERO/ALGEBRA}{Review sample-size and degenerate cases,
the zero-residual event argument, and the tagged quadratic coefficient.}

\item \textbf{Appendix C: supplementary experiments.} The original
distribution tables, heavy-tail comparisons, population-oracle
approximations, local-union comparison, and dimension-scaling study remain.
The infinite-variance cases are explicitly outside the stated assumptions.
The following nine added studies extend this evidence.
\item \textbf{E04: Dependence-strength sweep (\msapp{app:dependence-sensitivity}).}
Separates sensitivity to correlation from the single dependent setting in the body. Varying correlation at fixed marginal scales shows stable TSCP coverage and decreasing volume as dependence strengthens.

\item \textbf{E05: Coordinate-heterogeneity sweep (\msapp{app:heterogeneity-sensitivity}).}
Identifies when coordinate rescaling is useful by varying the largest-to-smallest noise-scale ratio. The common unscaled threshold becomes less efficient as heterogeneity increases, while remaining competitive at equal scales.

\item \textbf{E06: Target-coverage sweep (\msapp{app:alpha-small-sample}).}
Checks that the comparison is not specific to a 90 percent target. Targets of 80, 90, and 95 percent show the coverage-volume trade-off and the relative behavior of the methods.

\item \textbf{E07: Exchangeable contamination (\msapp{app:contamination}).}
Tests the sensitivity of mean and standard deviation scaling to outliers. Increasing row-wise contamination shows that near-nominal coverage can coexist with severe volume inflation.

\item \textbf{E08: CQR base-interval sensitivity (\msapp{app:cqr-sensitivity}).}
Checks dependence on the width of the fitted starting interval. The base-miscoverage sweep reveals CQHR's sensitivity to fitted widths and puts the main CQR comparison in context.

\item \textbf{E09: Shifted CQR scores (\msapp{app:cqr-sensitivity}).}
Examines additive shifts as an alternative to capping signed scores. Constants 100, 300, and 1000 give nearly identical recorded rectangles, equal to shifted TSCP-GWC within each run; this is not a general invariance theorem.

\item \textbf{E10: Rectangular shape-template baseline (\msapp{app:shape-template}).}
Adds the requested comparison with the Tumu et al. family. The two-dimensional KDE/mean-shift study compares coverage, normalized volume, and runtime for a specific rectangular template implementation, not every multimodal design.

\item \textbf{E11: Real-data coordinate diagnostics (\msapp{app:real-coordinate-audit}).}
Goes beyond the original two-point energy illustration with all 51 target coordinates. Paired length ratios and marginal coverages show where volume gains do and do not shorten individual intervals, including unfavorable student-data results.

\item \textbf{E12: Real-data search and fallback diagnostics (\msapp{app:real-data-diagnostics}).}
Measures actual branch use, inspected candidates, and construction times. No backward search or fallback occurs at miscoverage 0.1 in the tested splits; larger miscoverage levels exercise both branches without demonstrating a worst-case scan.

\item \textbf{Figures, references, and final checks.} The current draft uses
17 added experimental PDF assets and eight restyled original experimental
assets; several assets can belong to one study. The Tumu et al.\ reference
and both schematic annotations are retained. No numerical data or figure
PDF was changed in the final copyedit.
\checkauthor{DATA/SHIFT/REPRO/COLOR/METADATA}{Resolve sampling provenance,
the shift explanation/selection rule, run documentation, revision
highlighting, and template metadata before submission.}
\end{itemize}
\end{document}
